\section*{第 \ref{ch:g10:sport-and-health} 章——身体活动与健康}

\begin{solution}{exo:g10:sport-and-health:1}
心脏更大，每搏输出量更多，静息心率更低；
\omterm{def:g10:muscles-and-joints:muscle}{肌纤维}中的毛细血管和\omterm{def:g10:cells-common-unit:organelle}{线粒体}更多；红
\omterm{def:g10:cells-common-unit:cell}{细胞}更多；\omterm{def:g10:exercise-and-energy:glycogen}{糖原}储备更大（并且血糖
和血压的调节更好）。
\end{solution}

\begin{solution}{exo:g10:sport-and-health:2}
每搏输出量已经增大，因此同样的 \qty{5}{L/min} 只需
较少搏动即可输出，$Q = f \times V_s$ 中 $V_s$ 增大则需要更小的
$f$。
\end{solution}

\begin{solution}{exo:g10:sport-and-health:3}
每天约一小时中等强度的活动。不必如此：步行、骑车
上学、爬楼梯和做家务都算数——关键在于能提高\omterm{def:g10:exercise-and-energy:expenditure}{能量消耗}的
\omterm{def:g10:sport-and-health:training}{身体活动}。
\end{solution}

\begin{solution}{exo:g10:sport-and-health:4}
借助某些物质或方法提高成绩，即通过越过
身体的调节而非建立适应来实现。EPO 越过
对红\omterm{def:g10:cells-common-unit:cell}{细胞}数量的调控；\omterm{prop:g10:sport-and-health:doping}{兴奋剂}越过疲劳这一
保护心脏的预警信号。
\end{solution}

\begin{solution}{exo:g10:sport-and-health:5}
持续疲劳、静息心率上升、成绩下降
（还包括睡眠紊乱、反复感染、\omterm{def:g10:muscles-and-joints:muscle}{肌腱}痛或骨痛）。
\end{solution}

\begin{solution}{exo:g10:sport-and-health:6}
12 周后：\omterm{def:g10:sport-and-health:training}{训练}四次的组降低 $72 - 60 = 12$ 次，
\omterm{def:g10:sport-and-health:training}{训练}两次的组降低 $72 - 64 = 8$ 次。适应与需求成正比——
\omterm{def:g10:sport-and-health:training}{训练}越多，变化越大；
两组中，适应都随着接近新状态而放缓。
\end{solution}

\begin{solution}{exo:g10:sport-and-health:7}
$V_s = Q/f$：从 $5000/72 \approx \qty{69}{mL}$ 增至 $5000/54 \approx
\qty{93}{mL}$，增大为原来的 $72/54 \approx 1.33$ 倍。
\end{solution}

\begin{solution}{exo:g10:sport-and-health:8}
$1 - 0.55 = 45\%$，降低了。第一步，从最少活动到低活动
（1.00 到 0.80）带来的收益最大。
\end{solution}

\begin{solution}{exo:g10:sport-and-health:9}
\qty{2.0}{kg}，约占她质量的 $2.0/55 \approx 3.6\%$。她的血容量
下降，心率升高以维持输出量，体温上升，
成绩下降，抽筋的可能性增大。
\end{solution}

\begin{solution}{exo:g10:sport-and-health:10}
$10 \times 1.1^7 \approx \qty{19.5}{km}$，因此 \qty{20}{km} 要到
第八周才能达到：约两个月。
\end{solution}

\begin{solution}{exo:g10:sport-and-health:11}
每根\omterm{def:g10:muscles-and-joints:muscle}{肌纤维}通过增加肌原纤维而增粗，神经系统
学会同时募集更多\omterm{def:g10:muscles-and-joints:muscle}{肌纤维}并使之更同步地收缩；两者都会提高
力量。\omterm{def:g10:muscles-and-joints:muscle}{肌纤维}数目在成年后即固定不变。
\end{solution}

\begin{solution}{exo:g10:sport-and-health:12}
每升血液的氧含量增加 $56/45 \approx 1.24$ 倍，约 24\%；在
\omterm{def:g10:heart-lungs-effort:output}{心输出量}和氧提取不变时，$\dot V\!\mathrm{O_2}$max 至多按
同样的比例增加。但含 56\% \omterm{def:g10:cells-common-unit:cell}{细胞}的血液稠得多；
若心跳因睡眠而减慢，它就可能在冠状动脉或脑
血管中形成血栓。
\end{solution}

\begin{solution}{exo:g10:sport-and-health:13}
在高原上，身体会在真正缺氧时自行提高
EPO 的分泌量，幅度受控，缺氧消失时又会降低
这就是调节在起作用。注射 EPO 则强加了一个
调节本身绝不会选择的数值，既无缺氧作为依据，也无
制动机制。
\end{solution}

\begin{solution}{exo:g10:sport-and-health:14}
工作中的肌肉无需胰岛素即可摄取葡萄糖，而\omterm{def:g10:sport-and-health:training}{训练}过的肌肉
对葡萄糖更敏感；因此步行可直接降低血糖，
并减少每餐所需的胰岛素。药物可以降低血糖，
但没有消除病因——即不敏感、废用的肌肉——
依然如故，对心脏、骨骼和肌肉质量等步行也能
改善的方面毫无作用。
\end{solution}

\begin{solution}{exo:g10:sport-and-health:15}
适应是在休息期间逐步建立的，
即在反复、逐渐加大的负荷之间完成。一个月的高强度负荷
更可能导致过度\omterm{def:g10:sport-and-health:training}{训练}和损伤，而不是体能提高，因为\omterm{def:g10:muscles-and-joints:muscle}{肌腱}和骨骼
需要数月而非数周才能适应。而且，停止\omterm{def:g10:sport-and-health:training}{训练}后数周内，
收益便会消失：一次突击\omterm{def:g10:sport-and-health:training}{训练}无法为一整年“打好基础”。经常的中等强度
活动贯穿全年，能做到那一个月做不到的事。
\end{solution}

\begin{solution}{pb:g10:sport-and-health:1}
\textbf{1.} \omterm{def:g10:sport-and-health:training}{训练}前 $4800/72 \approx \qty{67}{mL}$；\omterm{def:g10:sport-and-health:training}{训练}后 $4800/56 \approx
\qty{86}{mL}$。

\textbf{2.} $42 \times 62 \approx \qty{2.6}{L/min}$；$50 \times 62 =
\qty{3.1}{L/min}$，增加了 19\%。

\textbf{3.} 最大输出量 $\dot V\!\mathrm{O_2}/0.150$：由
\num{17.4} 增至 \qty{20.7}{L/min}；心率 200 次/分时的最大每搏输出量：
由 \num{87} 增至 \qty{103}{mL}。

\textbf{4.} 心脏更大、更强，可解释每搏输出量；
若\omterm{def:g10:muscles-and-joints:muscle}{肌纤维}中的毛细血管和\omterm{def:g10:cells-common-unit:organelle}{线粒体}更多，
则可解释提取量的增加。

\textbf{5.} 26 周内增长为 4 倍，平均每周增长的倍数为 $4^{1/26}
\approx 1.055$，约每周 5.5\%，在 10\% 的规则之内（$1.1^{26}
\approx 12$ 则允许的增幅远大于此）。

\textbf{6.} 过度\omterm{def:g10:sport-and-health:training}{训练}：静息心率升高、睡眠受扰、
感染、成绩下降，修复缓慢的组织出现疼痛。

\textbf{7.} 胫骨的应力性损伤（
应力性骨折或骨表面炎症）：骨是所有组织中适应最慢的，
且负荷在两周内翻倍。

\textbf{8.} 两次\omterm{def:g10:sport-and-health:training}{训练}之间恢复不完全；身体
持续处于劳损状态，使得降低心率的适应
正在逆转。

\textbf{9.} 一周接近静养并保证睡眠（适应是在
休息中建立的）；随后两周距离减至此前的一半，不做速度\omterm{def:g10:sport-and-health:training}{训练}，
每周递增十分之一；胫骨须经医生诊治，并休息
至无痛为止（倾听疼痛的信号）；饮水、进食以促进修复。

\textbf{10.} 一种\omterm{prop:g10:sport-and-health:doping}{兴奋剂}：它掩盖了疲劳，这一警示保护
心脏免于力竭与过热。

\textbf{11.} 萨姆 $32 \times 80 \approx \qty{2.6}{L/min}$；亚历克斯 $48
\times 64 \approx \qty{3.1}{L/min}$。即便按绝对量计，亚历克斯的也更大；
日常活动——上楼梯、步行、承载自身体重——
牵动的是身体自身的质量，因此每千克体重对应的数值
才决定身体的感受。

\textbf{12.} 萨姆为其最大值的 $12/32 \approx 38\%$，亚历克斯为 $12/48 =
25\%$。萨姆的负荷超过其上限的三分之一，
以致气喘吁吁。

\textbf{13.} 不用的肌肉从血液中摄取的葡萄糖很少，
并对胰岛素失去敏感性；同样的饮食每年需要更多胰岛素，
血糖便开始升高。

\textbf{14.} 由 1.00 降至 0.70：风险约降低 30\%。

\textbf{15.} \omterm{def:g10:sport-and-health:training}{身体活动}是指任何提高\omterm{def:g10:exercise-and-energy:expenditure}{能量消耗}的运动，
如步行上班、走楼梯而非乘电梯、做家务
——每天累计半小时即可，无需体育项目或器材。

\textbf{16.} 按每年 1\% 计，10\% 的余量相当于十年的衰退。

\textbf{17.} 萨姆的峰值较低，且已经
流失骨量二十五年；他的腕部已越过
一次跌倒即会骨折的阈值，而亚历克斯的骨骼在
青少年期承受负荷而形成，此后持续维持，尚有十年余量。

\textbf{18.} 适应是靠使用来维持的结构，
而非永久获得的特性：需求停止时它会被拆除，
需求恢复时则重建。

\textbf{19.} 心脏与动脉疾病、2 型糖尿病、骨质疏松
（还包括部分癌症和抑郁症）。糖尿病：不活动的肌肉
对胰岛素失去敏感性，血糖升高，胰腺
负担过重。

\textbf{20.} 静息心率：54 对 78；$\dot V\!\mathrm{O_2}$max
\num{48} 对 \qty{32}{mL/min}（每千克体重）；每周约跑步三小时，
坚持了十四年。
\end{solution}
